\documentclass[11pt,a4paper]{article}

\usepackage[ngerman]{babel}
\usepackage[utf8]{inputenc}
\usepackage[T1]{fontenc}
\usepackage{geometry}
\usepackage{microtype}
\usepackage{tabularx}
\usepackage{array}
\usepackage{booktabs}
\usepackage{titlesec}
\usepackage{enumitem}
\usepackage{xcolor}
\usepackage{hyperref}
\usepackage{colortbl}

\geometry{a4paper, margin=2.5cm}

\definecolor{accent}{RGB}{40,60,90}

\titleformat{\section}{\normalfont\Large\bfseries\color{accent}}{\thesection}{1em}{}
\titleformat{\subsection}{\normalfont\large\bfseries\color{accent}}{\thesubsection}{1em}{}

\hypersetup{
    colorlinks=true,
    linkcolor=accent,
    urlcolor=accent,
    pdftitle={VR Climbing Speedrun - Projektdokumentation}
}

\setlist[itemize]{leftmargin=1.5em, itemsep=0.2em}

\newcommand{\arbeitszeile}[2]{#1 & #2 \\}

\begin{document}

% ---------------- Titelseite ----------------
\begin{titlepage}
    \centering
    \vspace*{2cm}

    {\Huge\bfseries\color{accent} VR Climbing Speedrun}\\[0.5em]
    {\Large Projektdokumentation}\\[2cm]

    {\large Virtual Reality: Konzeption und Anwendung immersiver Technologien}\\[3cm]

    \renewcommand{\arraystretch}{1.4}
    \begin{tabularx}{0.6\textwidth}{X c}
        \textbf{Person} & \textbf{Anteil} \\
        \arrayrulecolor{accent!40}\specialrule{0.4pt}{2pt}{4pt}
        Alexander Zitzmann & 25\,\% \\
        Kürsat Ayverdi     & 25\,\% \\
        Philipp Höpting    & 25\,\% \\
        Tim Fischer        & 25\,\% \\
    \end{tabularx}

    \vfill
    {\large \today}
\end{titlepage}

\tableofcontents
\newpage

% ---------------- 1. Projektbeschreibung ----------------
\section{Projektbeschreibung}

Das Projekt ist eine VR-Anwendung, bei der der Spieler durch physikbasiertes Klettern eine vertikale Spielwelt durchquert. Ziel ist es, die Strecke so schnell wie möglich zu absolvieren.
\\
Der Fokus liegt auf immersivem VR-Gameplay, direkter Handinteraktion und einem einfachen, wiederholbaren Speedrun-Loop.

% ---------------- 2. Vorgehensweise ----------------
\section{Vorgehensweise / Entwicklung}

Die Entwicklung erfolgte iterativ in mehreren Schritten:

\subsection{Prototyping Phase}

Zuerst wurde ein grundlegendes VR-Setup in Unity erstellt:

\begin{itemize}
    \item XR Origin Setup
    \item Handcontroller Integration
    \item Grundlegende Bewegung im Raum
\end{itemize}

\subsection{Entwicklung des Klettersystems}

Anschließend wurde das zentrale Gameplay-System entwickelt:

\begin{itemize}
    \item Klettern über Grip-Input
    \item Bewegung basiert auf Handpositionsdifferenz: Charakter bewegt sich entgegengesetzt zur Handbewegung
\end{itemize}

Ziel war es, ein natürliches und physiknahes Bewegungsgefühl zu erzeugen.

\subsection{Game Loop Implementierung}

Danach wurde der komplette Spielablauf aufgebaut:

\begin{itemize}
    \item Countdown (Ready -- Set -- Go)
    \item Start des Timers
    \item Aktivphase (Climbing)
    \item Zielerkennung über Trigger
    \item Stop des Timers, Restart der Szene
\end{itemize}

\subsection{UI Integration}

Zur Verbesserung der Spielerfahrung wurde ein UI-System integriert:

\begin{itemize}
    \item Timer-Anzeige
    \item Countdown-Anzeige
    \item Finish-Anzeige
    \item Scoreboard im Hauptmenü
\end{itemize}

\subsection{Audio Integration}

Ein einfacher Audio-Loop wurde implementiert:

\begin{itemize}
    \item Musik startet beim Run-Start
    \item Musik läuft während des gesamten Runs
\end{itemize}

% ---------------- 3. Funktionsweise ----------------
\section{Funktionsweise des Spiels}

\subsection{Startphase}

Beim Start des Spiels erscheint ein Countdown:

\begin{enumerate}
    \item ``Ready''
    \item ``Set''
    \item ``Go''
\end{enumerate}

Nach ``Go'' beginnt der Timer und der Spieler kann sich frei bewegen.

\subsection{Gameplay Phase}

Der Spieler bewegt sich durch:

\begin{itemize}
    \item Greifen von Kletterflächen
    \item Ziehen des Körpers durch Handbewegung
    \item Nutzung der Umgebung zur Fortbewegung
\end{itemize}

Gravitation wirkt permanent und beeinflusst die Bewegung.

\subsection{Zielphase}

Sobald der Spieler das Ziel erreicht:

\begin{itemize}
    \item Timer wird gestoppt
    \item Endzeit wird angezeigt
    \item Spiel bleibt kurz stehen
\end{itemize}

Danach erfolgt ein automatischer Restart.

% ---------------- 4. Technische Umsetzung ----------------
\section{Technische Umsetzung}

\subsection{Engine \& Frameworks}

\begin{itemize}
    \item Unity Engine
    \item XR Interaction Toolkit
    \item TextMeshPro für UI
\end{itemize}

\subsection{Hauptsysteme}

\paragraph{GameManager}
\begin{itemize}
    \item Steuert Timer
    \item Kontrolliert Spielzustände
    \item Startet und beendet Runs
\end{itemize}

\paragraph{FinishDetection}
\begin{itemize}
    \item Erkennt Zielkontakt über Trigger
    \item Stoppt den Run
\end{itemize}

\paragraph{Movement System}
\begin{itemize}
    \item VR-Klettern über Controller-Input
    \item Bewegung basiert auf Handdifferenz
    \item CharacterController für Physik
\end{itemize}

\paragraph{Scene Management}
\begin{itemize}
    \item Neustart des Spiels nach Finish
    \item Reset des kompletten Game States
\end{itemize}

% ---------------- 5. Spielarchitektur ----------------
\section{Spielarchitektur}

Das Projekt folgt einer klaren Trennung:

\begin{center}
    \large
    Input $\rightarrow$ Movement System $\rightarrow$ GameManager $\rightarrow$ UI / State Control $\rightarrow$ Scene Restart
\end{center}

% ---------------- 6. Entwicklungsprobleme ----------------
\section{Entwicklungsprobleme}

\subsection{Kantenprobleme}

Spieler konnten teilweise an Geometrie hängen bleiben.

\begin{quote}
    $\rightarrow$ \textbf{Lösung:} Optimierung der Collider-Struktur
\end{quote}

\subsection{Bewegungsgefühl}

Feintuning der Klettergeschwindigkeit war notwendig, um ein natürliches Gefühl zu erzeugen.

\subsection{Synchronisation UI / Timer}

Timer und UI mussten zentral über den GameManager synchronisiert werden.

% ---------------- 7. Ergebnis ----------------
\section{Ergebnis}

Das Endprodukt ist ein funktionaler VR-Speedrun-Kletterprototyp mit:

\begin{itemize}
    \item stabiler Klettermechanik
    \item funktionierendem Timer-System
    \item klarer Spielstruktur
    \item wiederholbarem Gameplay-Loop
\end{itemize}

% ---------------- 8. Fazit ----------------
\section{Fazit}

Das Projekt zeigt die erfolgreiche Umsetzung eines immersiven VR-Klettersystems mit Fokus auf einfache, intuitive Mechaniken und einem klaren Speedrun-Gameplay.

\end{document}
